\documentclass[10pt,a4paper]{extarticle}

\usepackage[a4paper, top=16mm, bottom=16mm, left=16mm, right=16mm]{geometry}
\usepackage[T1]{fontenc}
\usepackage{lmodern}
\usepackage{microtype}
\usepackage[table]{xcolor}
\usepackage{parskip}
\usepackage{enumitem}
\usepackage{titlesec}
\usepackage{mdframed}
\usepackage{multicol}
\usepackage{fancyhdr}
\usepackage{tabularx}
\usepackage{array}
\usepackage{hyperref}

% ── Colours ───────────────────────────────────────────────────────────────────
\definecolor{DG}{HTML}{1B5E20}
\definecolor{MG}{HTML}{2E7D32}
\definecolor{LG}{HTML}{43A047}
\definecolor{PG}{HTML}{F1F8E9}
\definecolor{CG}{HTML}{C8E6C9}
\definecolor{TEA}{HTML}{00696E}
\definecolor{TEAL}{HTML}{E0F2F1}
\definecolor{AMB}{HTML}{E65C00}
\definecolor{ALT}{HTML}{FFF3E0}
\definecolor{RED}{HTML}{C62828}
\definecolor{RLT}{HTML}{FFEBEE}
\definecolor{SLT}{HTML}{37474F}
\definecolor{INK}{HTML}{212121}
\definecolor{MUT}{HTML}{78909C}
\definecolor{PURP}{HTML}{4A148C}
\definecolor{PLLT}{HTML}{F3E5F5}

% ── Layout ────────────────────────────────────────────────────────────────────
\setlength{\parskip}{3pt}
\setlist[itemize]{leftmargin=1.4em, itemsep=1pt, parsep=0pt, topsep=2pt,
                  label=\textcolor{MG}{\textbullet}}

\titleformat{\section}
  {\bfseries\normalsize\color{white}}
  {}
  {0em}
  {\setlength\fboxsep{4pt}\colorbox{DG}}
\titlespacing{\section}{0pt}{10pt}{4pt}

% ── Header / Footer ───────────────────────────────────────────────────────────
\pagestyle{fancy}
\fancyhf{}
\renewcommand{\headrulewidth}{0pt}
\fancyhead[L]{\small\color{MUT}300,000 Streets --- Sample Q\&A}
\fancyhead[R]{\small\color{MUT}Regen Melbourne \texttimes{} RMIT \textbullet{} Hishikesh Phukan}
\fancyfoot[C]{\small\color{MUT}\thepage}

% ── Q&A environment ───────────────────────────────────────────────────────────
% Each Q&A pair: \qa{audience-tag}{tag-colour}{Question}{Answer}
\newcommand{\qa}[4]{%
  \begin{mdframed}[
    backgroundcolor=PG,
    linecolor=MG,
    linewidth=1.2pt,
    leftline=true, rightline=false, topline=false, bottomline=false,
    innerleftmargin=10pt, innerrightmargin=8pt,
    innertopmargin=6pt, innerbottommargin=7pt,
    skipabove=5pt, skipbelow=2pt
  ]
  \noindent
  \colorbox{#2}{\tiny\color{white}\textbf{#1}}\hspace{4pt}%
  {\bfseries\color{DG} Q:} {\bfseries\color{INK}#3}\\[3pt]
  {\bfseries\color{MG} A:} {\color{SLT}#4}
  \end{mdframed}%
}

% Variant for tricky / watch-out questions
\newcommand{\qatrap}[3]{%
  \begin{mdframed}[
    backgroundcolor=RLT,
    linecolor=RED,
    linewidth=1.2pt,
    leftline=true, rightline=false, topline=false, bottomline=false,
    innerleftmargin=10pt, innerrightmargin=8pt,
    innertopmargin=6pt, innerbottommargin=7pt,
    skipabove=5pt, skipbelow=2pt
  ]
  \noindent
  \colorbox{RED}{\tiny\color{white}\textbf{WATCH OUT}}\hspace{4pt}%
  {\bfseries\color{RED} Q:} {\bfseries\color{INK}#2}\\[3pt]
  {\bfseries\color{RED} A:} {\color{SLT}#3}
  \end{mdframed}%
}

\newcolumntype{L}[1]{>{\raggedright\arraybackslash}p{#1}}

% =============================================================================
\begin{document}

% ── Cover strip ───────────────────────────────────────────────────────────────
\colorbox{DG}{\parbox{\linewidth}{%
  \vspace{5pt}\hspace{8pt}%
  {\color{white}\large\bfseries 300,000 Streets of Melbourne --- Sample Q\&A}\\[2pt]
  \hspace{8pt}{\color{LG}\small Use this to practise before your presentation.
  Questions are grouped by audience type.}\\[2pt]
  \hspace{8pt}{\color{white!70!DG}\small
    \colorbox{TEA}{\tiny\color{white}\textbf{GENERAL}}\ non-technical audience\ \textbullet\
    \colorbox{MG}{\tiny\color{white}\textbf{TECHNICAL}}\ data science / engineering\ \textbullet\
    \colorbox{AMB}{\tiny\color{white}\textbf{POLICY}}\ council / planning\ \textbullet\
    \colorbox{RED}{\tiny\color{white}\textbf{WATCH OUT}}\ common traps}
  \vspace{5pt}
}}

\vspace{6pt}

% =============================================================================
\section{Results \& Findings}

\qa{GENERAL}{TEA}
{Preston High has the highest overall score --- why is it the most dangerous?}
{A great question, and the most important finding in this project. The overall score is an average of ten indicators. Preston scores very well on nine of them --- footpaths, air quality, things to do, and so on. But one indicator, HS2 ``easy to cross'', scored 0.4 out of 10. That means there is no safe pedestrian crossing at the school gate on a high-traffic arterial road. Under the Healthy Streets framework, a score that low on a safety-critical indicator immediately triggers Major severity, no matter how good everything else is. The average hides the crisis. That is precisely why you cannot rely on a single headline number.}

\qa{GENERAL}{TEA}
{What does HS2~=~0.4 actually look like on the ground?}
{Every school morning and afternoon, students at Preston High School walk up to Cooma Street --- a busy arterial road --- and there is no marked crossing, no traffic signals, no pedestrian island. They wait for a gap in traffic and walk across informally. That happens every single day. The 0.4 score reflects the near-total absence of safe crossing infrastructure at that specific point.}

\qa{GENERAL}{TEA}
{What is the 192\% crash increase --- is that just because more people are cycling?}
{Partly, yes --- active travel has grown post-COVID. But the 192\% figure is pedestrian and cyclist crashes combined, and the trend is steeper in Darebin than state-wide averages. More importantly, at William Ruthven Secondary College, 100\% of crashes near the school gate happen during arrival and pickup windows. That pattern points to structural infrastructure gaps at school times, not just more road users.}

\qa{POLICY}{AMB}
{Which school needs fixing first?}
{Preston High School, clearly. It is the only Major-severity school, and the root cause is a single, well-defined, costed fix: a signalised pedestrian crossing. Estimated cost is \$80,000--\$200,000. That is relatively modest infrastructure spend to resolve a daily safety crisis for hundreds of students. Reservoir High School is the second priority because it sits in the most disadvantaged catchment and has the most indicators below the safe threshold.}

% =============================================================================
\section{Framework \& Methodology}

\qa{GENERAL}{TEA}
{What is the Healthy Streets framework and why did you use it?}
{Healthy Streets was developed by Dr Lucy Saunders in 2015 and has been adopted by Transport for London as a formal planning standard. It defines ten measurable indicators of how pleasant and safe a street is for walking. We chose it because it is evidence-based, it has been validated across multiple cities, and it is legible to councils --- a score of 6.0 or higher means ``good'', below 6.0 means there is a problem. It translates data science output into something a planning officer can act on directly.}

\qa{POLICY}{AMB}
{Is 6.0 an arbitrary threshold?}
{It is not arbitrary --- 6.0 is the threshold defined in the original Healthy Streets framework documentation by Saunders (2015) and used operationally by Transport for London. It represents the point at which an indicator is considered to provide an acceptable pedestrian environment. Below 6.0 means the street fails that standard. We applied it as-is rather than redefining it, which makes our results comparable to any other city using the same framework.}

\qa{TECHNICAL}{MG}
{How do you translate raw field observations into a 0--10 score?}
{Each indicator has its own dedicated Python module in \texttt{src/scoring/}. For example, HS2 looks at whether a formal crossing exists within a defined distance of the school gate, what type it is (signalised, zebra, or none), and its estimated visibility. These are combined using a weighted scoring function with hard rules at the boundaries --- no crossing at all maps to a score near zero regardless of other factors. The exact functions are documented in the codebase and are fully reproducible.}

% =============================================================================
\section{Machine Learning}

\qa{TECHNICAL}{MG}
{Why Ridge regression and not something more powerful like XGBoost or a neural network?}
{Two reasons. First, the dataset has only three schools at scoring time. Any ensemble method or neural network would massively overfit with n=3. Ridge regression's regularisation is specifically designed to handle small samples and correlated features --- and OSM spatial features are highly collinear (shade ratio and tree count both describe the same physical reality). Second, interpretability matters here. A council officer asking ``why is this school flagged?'' deserves an answer in plain English, not a black-box prediction. Ridge coefficients are directly readable.}

\qa{TECHNICAL}{MG}
{How can you trust a model trained on only 3 data points?}
{We do not over-claim on the ML results. The model's primary purpose is to identify which indicators can be predicted from open data alone --- not to replace field surveys. With Leave-One-Out cross-validation, we get three independent test predictions, one per school. The mean MAE of 2.88 is honest about the uncertainty. The key insight --- that HS7 and HS10 can be automated while HS2 cannot --- is robust even with this small sample, because the difference in prediction error between those indicators is large and directional. We also note this limitation explicitly in the analysis.}

\qa{TECHNICAL}{MG}
{Why Leave-One-Out cross-validation?}
{With n=3, any train/test split would leave either zero or one sample in the test set. LOO-CV is the only statistically sound evaluation strategy: each school is held out once as the test set while the other two train the model. This gives us three unbiased error estimates. It is a well-established approach for small-sample evaluation and is used widely in medical and environmental research for exactly this reason.}

\qa{TECHNICAL}{MG}
{What does MAE~2.88 mean in practice --- is that good?}
{On a 0--10 scale, a mean absolute error of 2.88 means our predictions are off by roughly 3 points on average. That is not precise enough to replace field surveys. But it is sufficient for the intended use case: triage. If a model predicts a school's HS7 as 5.2 and the true value is 7.0, you still know it is worth a closer look. For the two automated indicators (HS7, HS10), the individual MAEs are lower. We use the model to flag schools for follow-up, not to replace the measurement.}

\qa{TECHNICAL}{MG}
{Why can HS7 and HS10 be automated but not HS2?}
{HS7 (feel safe) correlates strongly with the Crime Statistics Agency crime rate per 100,000 residents, which is a published, SA2-level open dataset. HS10 (clean air) maps directly to EPA AirWatch PM2.5 annual averages. Both are continuous, clean, government-maintained numbers that vary meaningfully between catchments. HS2 (easy to cross) cannot be automated because OpenStreetMap tags whether a crossing \emph{exists} --- not whether it is signalised, how visible it is, whether it has tactile pavers, or how fast the adjacent traffic moves. The physical quality of a crossing requires a human observer.}

% =============================================================================
\section{Equity \& Socio-economics}

\qatrap{WATCH OUT}
{Your r~=~0.84 correlation --- that is only three data points. That is not statistically significant, is it?}
{You are absolutely right, and we acknowledge this directly. With n=3, no Pearson correlation is statistically significant in a formal inferential sense. We do not claim it is. What we do say is that the relationship is directional and consistent with a substantial body of peer-reviewed literature --- Giles-Corti et al.\ (2016) in The Lancet found the same pattern across multiple cities. Our r=0.84 is \emph{consistent} with that evidence, not proof on its own. It motivates further study across all 200+ Melbourne schools, which is Phase~1 of the roadmap.}

\qa{POLICY}{AMB}
{What is SEIFA, and why does it matter here?}
{SEIFA stands for Socio-Economic Indexes for Areas, published by the Australian Bureau of Statistics using Census data. The specific index we use --- IRSD, Index of Relative Socio-economic Disadvantage --- ranks areas from decile~1 (most disadvantaged) to decile~10 (least). We joined it to each school's catchment at SA1 level using population-weighted averaging. It matters because lower SEIFA decile correlates with lower car ownership, which means more children are dependent on walking. If the streets those children walk are also the worst quality, that is a compounding disadvantage worth naming explicitly in a policy context.}

\qa{POLICY}{AMB}
{Is this an equity problem or just a geography problem --- maybe dangerous roads are everywhere, not just in poor areas?}
{Dangerous roads are everywhere, yes. But the pattern we found is that the most disadvantaged catchments in our sample have both the worst street safety scores AND the highest proportion of children who must walk rather than be driven. That combination --- worse streets, higher walking dependency --- is the equity problem. A dangerous road in an affluent suburb is still a problem, but it affects fewer children because most are driven. In Reservoir, lower household income means fewer cars, which means more children exposed to a footpath that is too narrow and a crossing 200 metres away.}

% =============================================================================
\section{Scalability \& Next Steps}

\qa{GENERAL}{TEA}
{You assessed three schools --- what about the other 300 in Melbourne?}
{That is exactly what Phase~1 of the roadmap addresses. Two of our ten indicators --- HS7 (feel safe) and HS10 (clean air) --- can be computed entirely from open government data with no field work. Running the automated pipeline on all 200-plus Melbourne secondary schools takes minutes. It produces a ranked list of schools by predicted risk, filtered by SEIFA disadvantage. Councils and the state government can then direct the limited pool of surveyors to the flagged schools first. Phase~2 would send surveyors to those schools for the remaining eight indicators. No code changes are required.}

\qa{POLICY}{AMB}
{How long would a full Melbourne triage actually take and what would it cost?}
{The automated Phase~1 --- running open-data indicators on all schools --- could be completed in days once data access is confirmed. A surveyor spends roughly two hours per school for the field-based indicators. At 200 schools, that is 400 surveyor-hours. Prioritising the highest-risk 50 schools first would cut that to 100 hours. The pipeline itself is open-source and free to run. The main cost is surveyor time.}

\qa{TECHNICAL}{MG}
{How do you keep the data current --- crash data changes every year?}
{The pipeline is designed for re-runs. VicRoads crash data is updated annually on data.vic.gov.au, and the download handler in \texttt{src/data/} fetches the latest file each time \texttt{run\_all.py} is called. OSM data is live via the Overpass API. SEIFA and Census data update on a five-year Census cycle. EPA and CSA data update annually. A scheduled re-run once a year would keep all outputs current. Post-intervention re-surveys would also retrain the Ridge model as the dataset grows.}

% =============================================================================
\section{Dashboard \& Technical Delivery}

\qa{GENERAL}{TEA}
{Why a website and not a report or spreadsheet?}
{A PDF report is static --- a council officer cannot filter recommendations by school, change a scenario, or zoom into a crash heatmap. The dashboard is interactive: you can pick a school, choose an intervention, and see the before/after instantly. It runs in any browser with no installation. We also export a QGIS-ready GeoPackage for council GIS teams who want to pull the data directly into their existing planning tools. The goal was to make the findings as easy to act on as possible, not just as easy to read.}

\qa{TECHNICAL}{MG}
{Why GitHub Pages and not a hosted backend?}
{Three reasons: zero running cost, zero maintenance, and zero security surface. All thirty scenario results are pre-computed and serialised to JSON at pipeline run time. The dashboard reads static JSON --- there is no server, no database, no API to maintain. For a proof-of-concept that needs to stay live for a year without any infrastructure budget, that is the right trade-off. If the project scales to real-time data feeds, a lightweight backend would be added at that point.}

\vspace{8pt}

% ── Bottom tip strip ──────────────────────────────────────────────────────────
\colorbox{SLT}{\parbox{\linewidth}{%
  \vspace{5pt}\hspace{8pt}%
  {\color{white}\bfseries\small Delivery tips for Q\&A}
  \vspace{3pt}
}}
\colorbox{SLT!15}{\parbox{\linewidth}{%
  \vspace{4pt}\hspace{8pt}%
  {\color{SLT}\small
  \textbf{Pause first.} Repeat the question back in one sentence before answering --- it buys thinking time and shows you understood. \quad
  \textbf{Quantify.} Vague answers lose marks. Always anchor to a number: 0.4, 192\%, r=0.84, \$80k, n=3. \quad
  \textbf{Own the limits.} Saying ``with only three schools this is directional, not inferential'' is a strength, not a weakness. Judges reward intellectual honesty.
  }
  \vspace{4pt}
}}

\end{document}
